\scanpage{72}
de $X_i$ suivant \unsure{son point|sa partie} à l'infini\ldots). Chaque $\hat X_i$ est automatiquement $\omega$-orienté, donc aussi
\[
\hat X=\coprod\hat X_i
\]

De cette façon, la catégorie envisagée (des surfaces topologiques ``spéciales'' $\omega$-orientées \inserted{-- avec \deleted{\unsure{leurs}}\unsure{homéomorphismes} $\omega$-orientés}) devient équivalente à celle des surfaces topologiques \inserted{$\omega$-orientées} à comp. conn. compactes \uline{munies} d'une partie \deleted{\unsure{f}} $S$ discrète, avec les isom. de couples $(\hat X,S)$.

De $\overline{\mathrm m}$, la catégorie des couples $\deleted{\unsure{(X,S)|(X,o,S)}}$ $(X,o,J)$ avec $(X,o)$ comme dessus ($o$ la $\omega$-orientation) et $J$ partie discrète de $X$, est équivalente à celle des systèmes $(\hat X,o,S,J)$, où $(\hat X,o)$ est une surface $\omega$-orientée à comp. conn compactes,
% archive sheet number: 71
